\documentclass[10pt,a4paper]{amsart}
\usepackage[margin=19mm]{geometry}
\usepackage{amsmath,amssymb,mathtools}
\usepackage[scr=rsfs,cal=euler]{mathalfa}
\usepackage{xfrac}
\usepackage[unicode,hidelinks,bookmarks=false]{hyperref}
\newcommand{\ie}{i.e.\ }
\newcommand{\cf}{cf.\ }
\newcommand{\resp}{respectively\ }
\newcommand{\Subsection}{\subsection}
\providecommand{\nobreakdash}{\nobreak\hbox{-}}
\setlength{\emergencystretch}{2em}
\multlinegap0pt
\usepackage{bm}
\usepackage{bbm}
\usepackage{dsfont}
\usepackage{xspace}
\usepackage{cancel}
\usepackage{mathtools}
\usepackage{amsthm}
\makeatletter
\@ifundefined{lem}{\newtheorem{lem}{Lemma}}{}
\@ifundefined{thm}{\newtheorem{thm}{Theorem}}{}
\makeatother
\title[Robust shape reconstruction of elastic impenetrable scattere]{Robust shape reconstruction of elastic impenetrable scatterers via monotonicity spectral sampling methods}
\author{Mengjiao Bai, Huaian Diao, Weisheng Zhou}
\date{Source: arXiv:2607.09062v1 (CC BY 4.0)}
\begin{document}
\maketitle

\noindent\emph{Source and licence.} Robust shape reconstruction of elastic impenetrable scatterers via monotonicity spectral sampling methods. Version arXiv:2607.09062v1 (CC BY 4.0), distributed under the Creative Commons Attribution 4.0 International licence, as recorded in the arXiv licence metadata for this exact version and shown on the abstract page https://arxiv.org/abs/2607.09062v1. Full rights notice: licenses/arxiv-2607.09062-CC-BY-4.0.txt.\par\smallskip

\noindent\textit{Editorial scope.} Counted unit: the Thm whose source label is \texttt{LocialWaveFounc} in the article (source0.tex lines 617--622, source label \texttt{LocialWaveFounc}), delivered together with the article's own complete original proof of it (source0.tex lines 624--716, one \texttt{proof} environment of 93 lines). No mathematics is written, repaired or re-derived: statement and proof are the article's own text line for line. Required context retained uncounted: eq:RestriH (lines 532--538); RangeHG (lines 542--555); Inquallity (lines 599--605); Dim (lines 607--613). Every definition, display and label of those lines is retained unchanged. Omitted from this package and not redistributed: 1--531, 539--541, 556--598, 606--606, 614--616, 623--623, 717--1798. Nothing inside a retained excerpt is omitted or altered: every retained body is delivered exactly as the article's lines are written, so no omission receipt of any kind accompanies this package. Citations: the retained excerpts carry no citation command at all, so no bibliography entry of the article is transcribed and no reference is redistributed. Reference adaptation: none was needed and none was made; every reference inside the delivered unit resolves inside it, so no unresolved reference or citation is produced. Printed slips kept literal: no printed word, symbol, hypothesis or number of the counted statement or of its proof was altered. Layout and class: the article's own class is \texttt{amsart} with packages including graphicx, amssymb, epstopdf, geometry, hyperref, booktabs, array, paralist, verbatim, subfig, tabularx, amsmath, amsfonts, amsthm; the wrapper is the lane's pinned \texttt{amsart} editorial wrapper at 10pt on A4, which restates the theorem-like environments the retained text needs and adds the article's own macro definitions that the retained text needs (none), with extra wrapper packages (bm, bbm, dsfont, xspace, cancel, mathtools). The delivered numbering is the unit's own. Assets: no retained span contains a figure environment, a graphics-inclusion command, a PDF-page inclusion, a table or a TikZ picture; no image file is copied into this package, the package declares no asset dependency and no image file is redistributed with it. Licence: version arXiv:2607.09062v1 of this article is distributed under the Creative Commons Attribution 4.0 International licence, as recorded in the arXiv licence metadata for this exact version and shown on the abstract page https://arxiv.org/abs/2607.09062v1. A full rights notice is delivered at licenses/arxiv-2607.09062-CC-BY-4.0.txt. Characters: every retained excerpt is pure ASCII, so the delivered engine, XeLaTeX with its default Latin Modern text font, reports no missing character.
\begin{equation}\label{eq:RestriH}
    \mathbf{H}_{\tau}^*
    =
    \mathbf{H}_B^* \mathbf{R}_{\tau}^*
    =
    \sqrt{8\pi\omega}\,\mathbf{G}_B^d \mathbf{S}_B \mathbf{R}_{\tau}^*.
\end{equation}

	\begin{thm}\label{RangeHG}
		Let $B, D \subset \mathbb R^2$ are open and bounded domain with Lipschitz  boundary, if $B \nsubseteq D$, let $\tau \subseteq \partial B \setminus \overline{D} $ is relatively open and $\mathbb R^2 \setminus \left( \overline{\tau \cup D }\right) $ is connected, then we have
			\begin{itemize}
				\item [(1).] When $D$ is a rigid impenetrable scatterer, it holds that
				\begin{align}\notag
					R\left( \mathbf H^*_{\tau} \right) \cap R\left( \mathbf G_D^d\right) =\left\lbrace 0\right\rbrace.
				\end{align}
				\item[(2).] When $D$ is a traction-free impenetrable scatterer, we obtain the following result
				\begin{align}\notag
					R\left( \mathbf H^*_{\tau} \right) \cap R\left( \mathbf G_D^n\right) =\left\lbrace 0\right\rbrace.
				\end{align}
			\end{itemize}
			
	\end{thm}

\begin{lem}\label{Inquallity}
	Assume $X, Y, Z$ be Hilbert spaces, let $A_1:X \to Y$ and $A_2: X \to Z$ are linear operators, then the following two statements are equivalent
	\begin{itemize}
		\item [(1).] There exists a constant $C>0$ such that $||A_1 x||_Y \leq C ||A_2 x||_Z, \quad \forall x \in X$.
		\item [(2).] $R \left( A^*_1 \right) \subseteq R \left( A^*_2 \right) $.
	\end{itemize}
\end{lem}

\begin{lem}\label{Dim}
	Assume $V, Z_1, Z_2 $ be subspaces of a vector space $Z$, if 
	$$
	Z_1 \cap Z_2= \left\lbrace 0 \right\rbrace \quad and \quad Z_1 \subseteq Z_2 +V,
	$$
	then $\rm {dim}\left(  Z_1 \right) \leqslant \rm {dim}\left(  V \right) $.
\end{lem}

\begin{thm}\label{LocialWaveFounc}
 Let $B, D  \subseteq \mathbb R^2$ be open and bounded domain with  Lipschitz boundary such that $\mathbb R^2 \backslash \overline {D}$ is connected. Assume $B \nsubseteq D$, then for any finite-dimensional subspace $V_1 \subseteq \left[L^{2}\left(\mathbb S\right) \right] ^2 $ , there exists a sequence $\left\lbrace  \mathbf f_n \right\rbrace  \subseteq  V_1^{\perp} $ such that
 $$
 \left\|\mathbf H_B \mathbf f_n\right\|_{\left[H^{1/2}\left(\partial B\right) \right] ^2 } \to \infty \quad and \quad \left\| \mathbf G_D^{d*} \mathbf f_n\right\|_{\left[H^{-1/2}\left(\partial D\right) \right] ^2 } \to 0, \qquad n\to \infty.
 $$
\end{thm}

\begin{proof}
Let $B, D  \subseteq \mathbb R^2$ is open and bounded domain with Lipschitz boundary such that $\mathbb R^2 \backslash \overline {D}$ is connected. Assume $B \nsubseteq D$, let $V_1 \subseteq \left[L^{2}\left(\mathbb S\right) \right] ^2 $ is a finite-dimensional subspace, therefore, the orthogonal projection onto $V_1$ is well-defined, and we denote it as $\mathbf P_1 :\left[L^{2}\left(\mathbb S\right) \right] ^2 \to \left[L^{2}\left(\mathbb S\right) \right] ^2$. 

Since $B \nsubseteq D $ and $\mathbb R^2 \backslash \overline {D}$ is connected, there exists a relatively open $\tau \subseteq \partial B\backslash \overline{D}$ such that $\mathbb R^2 \backslash \overline {\left( \tau \cup D\right) }$ is connected, from Theorem \ref{RangeHG} we have
\begin{align}\label{eq:RangeHtG}
R\left( \mathbf H^*_{\tau} \right) \cap R\left( \mathbf G_D^d\right) =\left\lbrace 0\right\rbrace .
\end{align}
We now turn our attention to \eqref{eq:RestriH}. Assume that $\omega^2$ is not a Dirichlet eigenvalue of the Navier operator $\Delta^*$ in $B$. Then both the single-layer potential operator $\mathbf{S}_B$ and the data-to-pattern operator $\mathbf{G}_B^d$ associated with $B$ are injective.

Moreover, the range of the extension operator $\mathbf R^*_{\tau} $ is infinite-dimensional, which implies that $ R \left( \mathbf H^*_{\tau} \right) = R \left( \sqrt{8\pi\omega}\mathbf{G}^d_B \mathbf S_B  \mathbf R^*_{\tau}   \right) $ is infinite-dimensional. Therefore, according to Lemma \ref{Dim} and \eqref{eq:RangeHtG} , we can obtain
$$
R \left( \mathbf H^*_{\tau} \right) \nsubseteq R\left( \mathbf G_D^d\right) +V_1=  R \left( \left[ \mathbf G_D^d \quad \mathbf P_1 \right] \right).
$$
Hence, utilizing Lemma \ref{Inquallity}, we can deduce that there does not exist a constant $C_1 > 0$ such that
\begin{equation}\begin{aligned}\label{eq:inver}
|| \mathbf H_{\tau} \mathbf g ||^2_{\left[H^{1/2}\left(\tau \right) \right] ^2 }
&\leq C_1^2 \left\| 
\begin{bmatrix}
	\mathbf G_D^{d*} \\ \mathbf P_1
\end{bmatrix}
\mathbf g \right\| ^2_{\left[H^{-1/2}\left(\partial D \right) \right] ^2  \times \left[L^{2}\left(\mathbb S \right) \right] ^2}\\
&=  C_1^2 \left( \left\| \mathbf G_D^{d*} \mathbf g \right\|^2_{\left[H^{-1/2}\left(\partial D \right) \right] ^2}+\left\| \mathbf P_1 \mathbf g \right\|^2_{\left[L^{2}\left(\mathbb S \right) \right] ^2} \right) .
\end{aligned}\end{equation}
Since $\mathbf P_1$ is an orthogonal projection operator, it follows that $\mathbf P_1$ is a self-adjoint operator, that is, $\mathbf P_1 = \mathbf P_1 ^*$. Therefore, from \eqref{eq:inver} for any $n \in \mathbb N$, there exists a vector valued function $\mathbf g_n \in \left[L^{2}\left(\mathbb S \right) \right] ^2  $ such that
$$
\left\| \mathbf H_{\tau} \mathbf g_n \right\|^2_{\left[H^{1/2}\left(\tau \right) \right] ^2 } > n^2 \left( \left\| \mathbf G_D^{d*} \mathbf g_n \right\|^2_{\left[H^{-1/2}\left(\partial D \right) \right] ^2}+\left\| \mathbf P_1 \mathbf g_n \right\|^2_{\left[L^{2}\left(\mathbb S \right) \right] ^2} \right).
$$
For any $n \in \mathbb N$, we denote $M_n:=\left\| \mathbf G_D^{d*} \mathbf g_n \right\|^2_{\left[H^{-1/2}\left(\partial D \right) \right] ^2}+\left\| \mathbf P_1 \mathbf g_n \right\|^2_{\left[L^{2}\left(\mathbb S \right) \right] ^2}$,  $\tilde{\mathbf g}_n:=\mathbf g_n /\left(\sqrt{n M_n} \right) $. On one hand, we have
\begin{equation}\notag 
\begin{aligned}
\left\| \mathbf H_{\tau} \tilde{\mathbf g}_n\right\|^2_{\left[H^{1/2}\left(\tau \right) \right] ^2 }
&= \frac{1}{n M_n} \left\| \mathbf H_{\tau} \mathbf g_n \right\|^2_{\left[H^{1/2}\left(\tau \right) \right] ^2 } \\
&>\frac{n}{M_n} \left( \left\|  \mathbf G_D^{d*} \mathbf g_n \right\|^2_{\left[H^{-1/2}\left(\partial D \right) \right] ^2}+\left\|  \mathbf P_1 \mathbf g_n \right\|^2_{\left[L^{2}\left(\mathbb S \right) \right] ^2} \right) \\
&= n,
\end{aligned}\end{equation}
which implies that 
$$
\left\| \mathbf H_{\tau} \tilde{\mathbf g}_n \right\|^2_{\left[H^{1/2}\left(\tau \right) \right] ^2 } > n \to \infty,\quad   n \to \infty .
$$

On the other hand, 
\begin{equation}\notag
\begin{aligned}
		\left\| \mathbf G_D^{d*} \tilde{\mathbf g}_n \right\|^2_{\left[H^{-1/2}\left(\partial D \right) \right] ^2}+\left\| \mathbf P_1 \tilde{\mathbf g}_n \right\|^2_{\left[L^{2}\left(\mathbb S \right) \right] ^2}
		&= \frac{1}{n M_n} \left(\left\| \mathbf G_D^{d*} \mathbf g_n \right\|^2_{\left[H^{-1/2}\left(\partial D \right) \right] ^2}+\left\| \mathbf P_1 \mathbf g_n \right\|^2_{\left[L^{2}\left(\mathbb S \right) \right] ^2} \right) \\
		&= \frac{1}{n}.
\end{aligned}\end{equation}
Thus, it yields that 
$$
\left\| \mathbf G_D^{d*} \tilde{\mathbf g}_n \right\|^2_{\left[H^{-1/2}\left(\partial D \right) \right] ^2}+\left\| \mathbf P_1 \tilde{\mathbf g}_n \right\|^2_{\left[L^{2}\left(\mathbb S \right) \right] ^2}= \frac{1}{n} \to 0, \quad n \to \infty .
$$
In summary, when $n \to \infty $, it can be deduced that
\begin{equation}\begin{aligned}\label{eq:HGPasymptotic}
  \left\| \mathbf H_{\tau} \tilde{\mathbf g}_n \right\|^2_{\left[H^{1/2}\left(\tau \right) \right] ^2 } \to \infty ,\quad \left\| \mathbf G_D^{d*} \tilde{\mathbf g}_n \right\|^2_{\left[H^{-1/2}\left(\partial D \right) \right] ^2} \to 0,\quad \left\| \mathbf P_1 \tilde{\mathbf g}_n \right\|^2_{\left[L^{2}\left(\mathbb S \right) \right] ^2} \to 0.
 \end{aligned}\end{equation}

For any $n \in \mathbb{N}$, we define $\mathbf{f}_n := \tilde{\mathbf{g}}_n - \mathbf{P}_1 \tilde{\mathbf{g}}_n$. By applying the triangle inequality, we obtain
\begin{equation}\begin{aligned}\label{eq:HInequality}
	    \left\| \mathbf H_{\tau} \mathbf f_n \right\|_{\left[H^{1/2}\left(\tau \right) \right]^2 }
		%&\geq \left| \left\| \mathbf H_{\tau} \tilde{\mathbf g}_n \right\|_{\left[H^{1/2}\left(\tau \right) \right]^2 } - \left\| \mathbf H_{\tau} \mathbf P_1 \tilde{\mathbf g}_n \right\|_{\left[H^{1/2}\left(\tau \right) \right]^2 } \right |\\
		&\geq \left\| \mathbf H_{\tau} \tilde{\mathbf g}_n \right\|_{\left[H^{1/2}\left(\tau \right) \right]^2 } - \left\| \mathbf H_{\tau} \mathbf P_1 \tilde{\mathbf g}_n \right\|_{\left[H^{1/2}\left(\tau \right) \right]^2 }\\
		&\geq \left\| \mathbf H_{\tau} \tilde{\mathbf g}_n \right\|_{\left[H^{1/2}\left(\tau \right) \right]^2 } - \left\| \mathbf H_{\tau} \right\|_{ \left[L^{2}\left(\mathbb S \right) \right]^2 \to \left[H^{1/2} \left(\tau \right) \right] ^2 } \left\|\mathbf P_1 \tilde{\mathbf g}_n \right\|_{\left[L^{2} \left(\mathbb S \right) \right] ^2}.
\end{aligned}\end{equation}
and
\begin{equation}\begin{aligned}\label{eq:GInequality}
		\left\| \mathbf G_D^{d*} \mathbf f_n \right\|_{\left[H^{-1/2}\left(\partial D \right) \right]^2 }
		&\leq \left\| \mathbf G_D^{d*} \tilde{\mathbf g }_n\right\|_{\left[H^{-1/2}\left(\partial D \right) \right]^2 }+ \left\| \mathbf G_D^{d*} \mathbf P_1 \tilde{\mathbf g }_n\right\|_{\left[H^{-1/2}\left(\partial D \right) \right]^2 } \\
		&\leq \left\| \mathbf G_D^{d*} \tilde{\mathbf g }_n\right\|_{\left[H^{-1/2}\left(\partial D \right) \right]^2 }+ \left\| \mathbf G_D^{d*}\right\|_{\left[L^{2} \left(\mathbb S \right) \right] ^2 \to \left[H^{-1/2}\left(\partial D \right) \right]^2} \left\|\mathbf P_1 \tilde{\mathbf g }_n\right\|_{\left[L^{2} \left(\mathbb S \right) \right] ^2}.
\end{aligned}\end{equation}

It follows from \eqref{eq:HGPasymptotic}, \eqref{eq:HInequality}, and \eqref{eq:GInequality} that 
$\|\mathbf{H}_{\tau} \mathbf{f}_n\|_{[H^{1/2}(\tau)]^2} \to \infty$ and 
$\|(\mathbf{G}_D^d)^* \mathbf{f}_n\|_{[H^{-1/2}(\partial D)]^2} \to 0$ as $n \to \infty$. 
Recalling the definition $\mathbf{H}_{\tau} = \mathbf{R}_{\tau} \mathbf{H}_B$, we have
\begin{equation}\notag
\begin{aligned}
\|\mathbf{H}_{\tau} \mathbf{f}_n\|_{[H^{1/2}(\tau)]^2} 
&= \|\mathbf{R}_{\tau} \mathbf{H}_B \mathbf{f}_n\|_{[H^{1/2}(\partial B)]^2} \\
& \leq \left\|\mathbf R_{\tau}\right\|_{\left[H^{1/2}\left(\partial B \right) \right]^2 \to \left[H^{1/2}\left(\tau \right) \right]^2 } \left\| \mathbf H_{B} \mathbf f_n \right\|_{\left[H^{1/2}\left(\partial B \right) \right]^2 },
\end{aligned}
\end{equation}
which implies
\begin{equation}\notag 
\|\mathbf{H}_B \mathbf{f}_n\|_{[H^{1/2}(\partial B)]^2} 
\geq \frac{\|\mathbf{H}_{\tau} \mathbf{f}_n\|_{[H^{1/2}(\tau)]^2}}{\left\|\mathbf R_{\tau}\right\|_{\left[H^{1/2}\left(\partial B \right) \right]^2 \to \left[H^{1/2}\left(\tau \right) \right]^2 }}.
\end{equation}
Since $\|\mathbf{H}_{\tau} \mathbf{f}_n\|_{[H^{1/2}(\tau)]^2} \to \infty$ as $n \to \infty$ and the norm of the restriction operator $\mathbf{R}_{\tau}$ is bounded, it follows that
\begin{equation}\notag
\|\mathbf{H}_B \mathbf{f}_n\|_{[H^{1/2}(\partial B)]^2} \to \infty \quad \text{as } n \to \infty.
\end{equation}
This, together with \eqref{eq:GInequality}, completes the proof.

\end{proof}

\end{document}
